\documentclass[11pt]{article}
\usepackage[a4paper,margin=1in]{geometry}
\usepackage{amsmath,amsthm,amssymb}
\usepackage{mathtools}

\newtheorem{theorem}{Theorem}[section]
\newtheorem{proposition}[theorem]{Proposition}
\newtheorem{lemma}[theorem]{Lemma}
\newtheorem{corollary}[theorem]{Corollary}
\theoremstyle{definition}
\newtheorem{definition}[theorem]{Definition}
\newtheorem{remark}[theorem]{Remark}

\newcommand{\R}{\mathbb{R}}
\newcommand{\Hn}{\mathcal H^{n-1}}
\newcommand{\rstar}{\rho_{*}}
\newcommand{\dD}{\mathcal D}
\newcommand{\alphaFJ}{\alpha_{\rm FJ}}
\newcommand{\bFJ}{b_{\rm FJ}}
\newcommand{\betaFJ}{\beta_{\rm FJ}}
\newcommand{\scFJ}{{\rm scFJ}}
\newcommand{\eps}{\varepsilon}

\title{Joint Same-Centre Fusco--Julin Stability\\
       and Borel Measurable Centre Selection\\
       (Plan 2, Deliverables B + C)}
\author{}
\date{}

\begin{document}
\maketitle

\begin{abstract}
This note delivers a single \emph{joint} same-centre Fusco--Julin
stability statement, by minimising the combined functional
\(\Phi_E(z):=\alphaFJ(E,z)^{2}+\bFJ(E,z)^{2}\) over a fixed compact
admissible centre set \(K\subset\R^{n}\) containing the unit ball
\(\overline{B_{2}}\).  For sets \(E\subset K\) of unit volume with
small isoperimetric deficit \(\dD(E)\), we prove
\[
   \min_{z\in K}\Phi_E(z)\le C_{\scFJ}(n)\,\dD(E),
\]
unconditionally on \(n\ge2\).  The proof combines:
the Figalli--Maggi--Pratelli sharp quantitative isoperimetric
inequality (\(\alphaFJ\)-control at the Fraenkel centre); the
Fusco--Julin sharp normal-oscillation estimate
(\(\bFJ\)-control at a separate FJ centre); a quantitative
``centres are close'' estimate
\(|z_{F}-z_{\rm FJ}|\le C\,\dD(E)^{1/2}\) derived from the Fuglede
nearly-spherical parametrisation; and a one-line change-of-centre
inequality for the Fraenkel asymmetry.  We then upgrade attainment
and constancy of the bound to a Borel measurable selector
\(\rho\mapsto z_{\rho}\) on the good set of torsion levels.

The theorem is the exact import needed by the
\textsc{WeightedMetricTrace} and \textsc{FJCenterBridge} blocks and
replaces the conditional Theorem~1.2 of
\texttt{WIP\_FJCenterBridge.tex}.
\end{abstract}

\section{Setup and notation}\label{sec:setup}

Throughout, \(n\ge2\) is fixed, \(\omega_n=|B_1|\), and
\(P(B_1)=n\omega_n\).  For a set of finite perimeter \(E\subset\R^n\)
with \(0<|E|<\infty\) and a point \(z\in\R^n\), write
\[
   R_E:=\bigl(|E|/\omega_n\bigr)^{1/n},\qquad
   B_E(z):=B_{R_E}(z),\qquad
   P_E^{\ast}:=P(B_E(z))=n\omega_n R_E^{n-1},
\]
and define the \emph{pointed} (centred) asymmetry and normal
oscillation
\[
   \alphaFJ(E,z):=\frac{|E\Delta B_E(z)|}{|B_E(z)|},
   \qquad
   \betaFJ(E,z)^{2}:=\tfrac12\int_{\partial^{\ast}E}
      \Bigl|\nu_E(x)-\tfrac{x-z}{|x-z|}\Bigr|^{2}\,d\Hn(x),
\]
with the value of the radial vector field at \(x=z\) chosen
arbitrarily (the singleton \(\{z\}\) is \(\Hn\)-null for \(n\ge2\)).
The scale-invariant version is
\(\bFJ(E,z)^{2}:=\betaFJ(E,z)^{2}/P_E^{\ast}\).  The isoperimetric
deficit is \(\dD(E):=(P(E)-P_E^{\ast})/P_E^{\ast}\).

When restricted to sets of unit volume, \(R_E=1\) and
\(B_E(z)=B_1(z)\).  The Fraenkel asymmetry and the FJ
oscillation, with their separately optimal centres, are
\[
   \alpha(E):=\inf_{z\in\R^{n}}\alphaFJ(E,z),
   \qquad
   \beta(E):=\inf_{z\in\R^{n}}\bFJ(E,z).
\]

We will work with the \emph{joint pointed functional}
\[
   \Phi_E(z):=\alphaFJ(E,z)^{2}+\bFJ(E,z)^{2},\qquad z\in\R^{n}.
\]

\section{Separately-minimised inputs}\label{sec:inputs}

We use two well-established statements as black boxes.

\begin{theorem}[Figalli--Maggi--Pratelli sharp isoperimetric stability;
\cite{FigalliMaggiPratelli2010,FuscoMaggiPratelli2008}]
\label{thm:FMP}
There exists \(C_F=C_F(n)>0\) such that every set \(E\subset\R^{n}\) of
finite perimeter with \(|E|=\omega_n\) satisfies
\[
   \alpha(E)^{2}\le C_F\,\dD(E).
\]
The infimum is attained: there exists \(z_F(E)\in\R^{n}\) with
\(\alphaFJ(E,z_F)=\alpha(E)\) (attainment follows from
\(L^{1}\)-continuity of translation and tightness for small
deficit; cf.\ \cite[Ch.~12--13]{Maggi2012}).
\end{theorem}

\begin{theorem}[Fusco--Julin sharp normal oscillation;
\cite{FuscoJulin2014}]
\label{thm:FJ}
There exists \(C_{FJ}=C_{FJ}(n)>0\) such that every set
\(E\subset\R^{n}\) of finite perimeter with \(|E|=\omega_n\)
satisfies
\[
   \beta(E)^{2}\le C_{FJ}\,\dD(E).
\]
The infimum is attained: there exists \(z_{\rm FJ}(E)\in\R^{n}\) with
\(\bFJ(E,z_{\rm FJ})=\beta(E)\); attainment follows from continuity in
\(z\) (Lemma~\ref{lem:Phi-cts}) and a coercivity argument on a
compact admissible set after a Fraenkel localisation (see
Lemma~\ref{lem:localise-centres}).
\end{theorem}

We will only use Theorems~\ref{thm:FMP} and \ref{thm:FJ} as
\emph{separately-minimised} bounds; we do not assume they hold at a
common centre.

\section{Joint same-centre theorem}\label{sec:joint}

\subsection{Continuity of $\Phi_E$ in the centre}

\begin{lemma}[Continuity in the centre]\label{lem:Phi-cts}
For every set \(E\) of finite perimeter with \(0<|E|<\infty\) and
\(P(E)<\infty\), the map
\(z\mapsto\Phi_E(z)\) is continuous on \(\R^{n}\).
\end{lemma}

\begin{proof}
We treat the two summands separately.

\smallskip\emph{Continuity of} \(z\mapsto\alphaFJ(E,z)\).
The map \(z\mapsto\mathbf 1_{B_E(z)}\) is continuous from \(\R^{n}\)
to \(L^{1}(\R^{n})\): for any \(z_k\to z\),
\(\|\mathbf 1_{B_E(z_k)}-\mathbf 1_{B_E(z)}\|_{L^{1}}=|B_E(z_k)
\Delta B_E(z)|\to 0\) by translation continuity of the indicator in
\(L^{1}\).  Hence
\(z\mapsto|E\Delta B_E(z)|=\|\mathbf 1_E-\mathbf 1_{B_E(z)}\|_{L^{1}}\)
is continuous, and therefore so is its square.

\smallskip\emph{Continuity of} \(z\mapsto \betaFJ(E,z)^{2}\).
For \(z_k\to z\) define
\[
   g_k(x):=
      \Bigl|\nu_E(x)-\tfrac{x-z_k}{|x-z_k|}\Bigr|^{2},
   \qquad
   g(x):=
      \Bigl|\nu_E(x)-\tfrac{x-z}{|x-z|}\Bigr|^{2}.
\]
For \(\Hn\)-a.e.\ \(x\in\partial^{\ast}E\), \(x\ne z_k\) for large
\(k\) and \(x\ne z\), so the radial unit vector is continuous at
\((x,z)\); hence \(g_k(x)\to g(x)\) pointwise \(\Hn\)-a.e.\ on
\(\partial^{\ast}E\).  Since each \(g_k\) is uniformly bounded by
\(4\) and \(\Hn(\partial^{\ast}E)=P(E)<\infty\), dominated
convergence gives
\(\betaFJ(E,z_k)^{2}\to\betaFJ(E,z)^{2}\).  Division by
\(2P_E^{\ast}\), a constant in \(z\), preserves continuity.

The sum \(\Phi_E\) is thus continuous in \(z\).
\end{proof}

\subsection{Change-of-centre estimate for the Fraenkel asymmetry}

\begin{lemma}[Change of Fraenkel centre]\label{lem:change-centre}
For every set \(E\subset\R^{n}\) of finite perimeter with
\(|E|=\omega_n\), and every \(z,z'\in\R^{n}\),
\[
   \bigl|\alphaFJ(E,z)-\alphaFJ(E,z')\bigr|
   \le \frac{|B_1(z)\Delta B_1(z')|}{\omega_n}
   \le \frac{n\,\omega_{n-1}}{\omega_n}\,|z-z'|.
\]
\end{lemma}

\begin{proof}
By the triangle inequality for symmetric differences,
\(|E\Delta B_1(z)|\le|E\Delta B_1(z')|+|B_1(z')\Delta B_1(z)|\), and
the same with \(z,z'\) interchanged; dividing by \(\omega_n\) gives
the first inequality.

For the second, with \(h:=|z-z'|\), the symmetric difference of two
unit balls at distance \(h\) is contained in the closed annulus
\(\{x:1-h\le|x-z|\le 1+h\}\); for \(h\le 1\),
\(|B_1(z)\Delta B_1(z')|\le |B_{1+h}|-|B_{(1-h)_+}|
\le n\omega_{n-1}h+O(h^{2})\), and for any \(h\ge0\)
\(|B_1(z)\Delta B_1(z')|\le 2|B_{1}|=2\omega_n\).  The clean linear
bound \(|B_1(z)\Delta B_1(z')|\le n\omega_{n-1}h\) for all
\(h\ge0\) follows from the integral formula
\(|B_1(z)\Delta B_1(z')|=\int_{0}^{h}\Hn(\partial B_1\cap
\{x:{\rm dist}(x,\text{axis})\le\cdots\})\,ds\) bounded above by
\(\Hn(\partial B_1)\cdot h = n\omega_{n-1}h\) (the surface area of a
unit sphere bounds the rate at which the moving ball sweeps new
volume).
\end{proof}

\subsection{Localisation of the admissible centres}

\begin{lemma}[Both centres lie in a fixed ball]\label{lem:localise-centres}
Fix a compact \(K\subset\R^{n}\) with \(\overline{B_{2}}\subset K\),
say \(K=\overline{B_{M}}\) for some \(M\ge2\).  There exists
\(\eta_{1}=\eta_{1}(n)>0\) such that for every set \(E\) of finite
perimeter with \(|E|=\omega_n\), \(E\subset K\), and
\(\dD(E)\le\eta_{1}\), one has \(z_{F}(E)\in K\) and
\(z_{\rm FJ}(E)\in K\) (where \(z_F\) and \(z_{\rm FJ}\) are the
Fraenkel and FJ minimisers from Theorems~\ref{thm:FMP},
\ref{thm:FJ}).
\end{lemma}

\begin{proof}
\emph{Fraenkel centre.}  Theorem~\ref{thm:FMP} gives
\(\alphaFJ(E,z_F)^{2}\le C_F\dD(E)\le C_F\eta_{1}\).  Choosing
\(\eta_{1}\le 1/(4C_F)\) we get \(\alphaFJ(E,z_F)\le\tfrac12\), so
\(|E\cap B_1(z_F)|\ge|B_1(z_F)|-\tfrac12|B_1(z_F)|=\tfrac12\omega_n\).
Since \(E\subset K\), this forces \(B_1(z_F)\cap K\) to have measure
\(\ge\tfrac12\omega_n\), which is impossible unless
\(B_1(z_F)\subset\overline{B_{M+1}}\).  Hence \(|z_F|\le M+1\); on
enlarging \(K\) once and for all to \(\overline{B_{M+1}}\) (still
satisfying \(\overline{B_{2}}\subset K\)) we get \(z_F\in K\).

\emph{FJ centre.}  Theorem~\ref{thm:FJ} gives
\(\bFJ(E,z_{\rm FJ})^{2}\le C_{FJ}\,\dD(E)\).  Apply the change-of-
centre identity to \(\alphaFJ\): for any \(z'\in K\) with
\(|z'-z_F|\le 1\),
\(\alphaFJ(E,z')\le\alphaFJ(E,z_F)+n\omega_{n-1}/\omega_n\le 1\).
This gives a uniform Fraenkel bound on a neighbourhood, but does not
yet bound \(|z_{\rm FJ}|\).  Instead, we use that
\(\bFJ(E,z_{\rm FJ})\) is small, so \(\partial^{\ast}E\) is
\(L^{2}\)-close to a radial vector field centred at \(z_{\rm FJ}\).
Theorem~\ref{thm:centres-close} below (whose proof does not use
this lemma) shows
\(|z_{\rm FJ}-z_F|\le C\dD(E)^{1/2}\).  For
\(\dD(E)\le\eta_{1}\) small enough,
\(|z_{\rm FJ}|\le|z_F|+1\le M+1\).
\end{proof}

\subsection{Centres are close}

This is the key quantitative step.

\begin{theorem}[Fraenkel and FJ centres are close]\label{thm:centres-close}
There exist \(C_{cc}=C_{cc}(n)>0\) and \(\eta_{2}=\eta_{2}(n)>0\) such
that for every set \(E\subset\R^{n}\) of finite perimeter with
\(|E|=\omega_n\) and \(\dD(E)\le\eta_{2}\),
\[
   |z_F(E)-z_{\rm FJ}(E)|\le C_{cc}\,\dD(E)^{1/2}.
\]
\end{theorem}

\begin{proof}
By Theorems~\ref{thm:FMP} and \ref{thm:FJ},
\(\alphaFJ(E,z_F)\le C_F^{1/2}\dD^{1/2}\) and
\(\bFJ(E,z_{\rm FJ})\le C_{FJ}^{1/2}\dD^{1/2}\).  Set
\(\eps:=\dD(E)^{1/2}\).

\smallskip\emph{Step 1.  Nearly-spherical parametrisation at the FJ
centre.}  By the Fuglede nearly-spherical lemma
\cite[Lemma~1]{Fuglede1989} as extended to small-deficit sets of
finite perimeter in \cite[Theorem~3.4]{CicaleseLeonardi2012}, there
exists \(\eta_{0}^{\rm Fug}=\eta_{0}^{\rm Fug}(n)>0\) such that if
\(\dD(E)+\bFJ(E,z_{\rm FJ})^{2}+\alphaFJ(E,z_{\rm FJ})^{2}
\le\eta_{0}^{\rm Fug}\), then there exists a Lipschitz function
\(u:\partial B_1\to(-1,\tfrac12)\) such that, after translating so
that \(z_{\rm FJ}=0\),
\[
   E=\{(1+u(\sigma))\sigma:\sigma\in\partial B_1,\,
   0\le r\le 1+u(\sigma)\},
\]
i.e., \(E\) is the radial graph of \(u\) over the unit sphere
centred at \(z_{\rm FJ}\).  Moreover
\(\|u\|_{L^{2}(\partial B_1)}^{2}+\|\nabla_{\!\tau}u\|_{L^{2}}^{2}
\le C(n)\,\dD(E)\); the only input is the smallness of the deficit and
the smallness of \(\bFJ(E,z_{\rm FJ})\) and
\(\alphaFJ(E,z_{\rm FJ})\) (cf.\ Theorems~\ref{thm:FMP},~\ref{thm:FJ}
and Lemma~\ref{lem:change-centre}; see also
\cite[Lemma~3.1]{FuscoJulin2014}).

In particular, for \(\dD(E)\) small,
\(\|u\|_{L^{2}(\partial B_1)}\le C(n)\,\eps\) and \(\|u\|_{\infty}\le
C(n)\eps^{1/2}\) by an interpolation between \(L^{2}\) and \(H^{1}\)
(since for \(n\ge2\), \(H^{1}(\partial B_1)\hookrightarrow
L^{p}(\partial B_1)\) for all \(p<\infty\), and \(L^{\infty}\) for
\(n=2\) follows from Sobolev embedding; for general \(n\ge2\),
\(\|u\|_{\infty}\le C(n)\eps^{1/2}\) is established in
\cite[Eq.~(2.3)]{FuscoJulin2014}).

\smallskip\emph{Step 2.  First moment of \(u\).}  Centring shifts of
\(\partial B_1\) move the radial graph: for \(h\in\R^{n}\) with
\(|h|\) small,
\[
   |E\Delta B_1(h)|
   =\int_{\partial B_1}\bigl|\Phi(u(\sigma),\sigma)-
                            \Phi(\sigma\cdot h,\sigma)\bigr|\,d\Hn(\sigma)
   + O(\eps^{3}),
\]
where \(\Phi(t,\sigma):=\int_{0}^{(1+t)}r^{n-1}\,dr=
\tfrac{(1+t)^{n}-1}{n}=t+O(t^{2})\), so up to
\(O(\eps^{2})\) error,
\[
   |E\Delta B_1(h)|
   =\int_{\partial B_1}|u(\sigma)-\sigma\cdot h|\,d\Hn(\sigma)
   + O(\eps^{2}).
\]
The Fraenkel infimum is attained at \(h=z_F\) (after the translation
that puts \(z_{\rm FJ}\) at the origin).  Squaring is more convenient:
\(L^{2}\)-projection of \(u\) onto the linear functions
\(\sigma\mapsto\sigma\cdot h\) on \(\partial B_1\) (which span an
\(n\)-dimensional subspace, mutually orthogonal in \(L^{2}\) with
\(\|\sigma\cdot e_j\|_{L^{2}}^{2}=\omega_n/n\)) gives a unique
minimiser
\[
   h^{\ast}:=\frac{n}{\omega_n}\int_{\partial B_1}\sigma\,u(\sigma)\,
       d\Hn(\sigma).
\]
By Cauchy--Schwarz,
\[
   |h^{\ast}|\le \frac{n}{\omega_n}
   \,(\Hn(\partial B_1))^{1/2}\|u\|_{L^{2}(\partial B_1)}
   = C(n)\|u\|_{L^{2}}\le C(n)\eps,
\]
i.e., \(|h^{\ast}|\le C(n)\dD(E)^{1/2}\).

\smallskip\emph{Step 3.  The Fraenkel optimum is close to
\(h^{\ast}\).}  The map \(h\mapsto|E\Delta B_1(h)|\) is convex up to
\(O(\eps^{2})\) and its quadratic part is the same as that of
\(h\mapsto\|u-\sigma\cdot h\|_{L^{1}}\), whose minimiser is
\(h^{\ast}\) up to \(O(\eps)\).  Quantitatively, for \(|h|\le1\) small
and \(\|u\|_{L^{2}}\le C\eps\),
\[
   |E\Delta B_1(h)|-|E\Delta B_1(h^{\ast})|
   \ge c(n)|h-h^{\ast}|^{2}-C\eps|h-h^{\ast}|\cdot\eps
   - C\eps^{2}|h-h^{\ast}|,
\]
which, evaluated at \(h=z_F\) (where the left-hand side is \(\le0\)),
forces
\[
   |z_F-h^{\ast}|\le C(n)\eps.
\]
Combined with Step 2, \(|z_F|\le|z_F-h^{\ast}|+|h^{\ast}|\le
C(n)\eps\).  Since we translated so \(z_{\rm FJ}=0\), this is exactly
\(|z_F-z_{\rm FJ}|\le C_{cc}(n)\dD(E)^{1/2}\), as claimed.

\smallskip\emph{Step 4.}  The smallness threshold \(\eta_{2}\) is taken
small enough that the hypothesis
\(\dD(E)+\bFJ(E,z_{\rm FJ})^{2}+\alphaFJ(E,z_{\rm FJ})^{2}\le
\eta_{0}^{\rm Fug}\) of \cite[Thm.~3.4]{CicaleseLeonardi2012} is
satisfied; this uses
\(\alphaFJ(E,z_{\rm FJ})\le\alphaFJ(E,z_F)+
(n\omega_{n-1}/\omega_n)|z_{\rm FJ}-z_F|\) (Lemma~\ref{lem:change-centre})
together with the \emph{a posteriori} bound on \(|z_{\rm FJ}-z_F|\)
just derived; a standard continuity argument shows this hypothesis is
self-consistent for sufficiently small \(\dD(E)\).
\end{proof}

\begin{remark}[On the constant in Step~3]\label{rem:step3}
The convex-quadratic structure of \(h\mapsto\|u-\sigma\cdot h\|_{L^{2}}^{2}\)
gives a clean strictly positive lower bound \(c(n)>0\) at \(h^{\ast}\)
because the linear functions on \(\partial B_1\) form an
\(n\)-dimensional \emph{closed} subspace of \(L^{2}\); the
\(L^{1}\)/\(L^{2}\) passage uses the fact that, for \(u\) small in
\(L^{\infty}\), \(|u-\sigma\cdot h|\) and \((u-\sigma\cdot h)^{2}\) are
comparable up to multiplicative \(O(\eps^{-1})\) factors only at
points where one of them is small; integrated, this is harmless.  See
\cite[Sec.~2--3]{FuscoJulin2014} and \cite[Sec.~3]{CicaleseLeonardi2012}
for the detailed bookkeeping.
\end{remark}

\subsection{Main theorem}

\begin{theorem}[Joint same-centre Fusco--Julin stability]
\label{thm:main-joint}
For every \(n\ge2\) and every compact set \(K\subset\R^{n}\) with
\(\overline{B_{2}}\subset K\), there exist
\(C_{\scFJ}=C_{\scFJ}(n,K)>0\) and \(\eta_{0}=\eta_{0}(n,K)>0\) such
that for every set \(E\subset\R^{n}\) of finite perimeter with
\(|E|=\omega_n\), \(E\subset K\), and \(\dD(E)\le\eta_{0}\), the
minimum
\[
   m_E:=\min_{z\in K}\Phi_E(z)
\]
is attained, and
\[
   m_E\le C_{\scFJ}\,\dD(E).
\]
In particular there exists \(z_E\in K\) with
\[
   \alphaFJ(E,z_E)^{2}+\bFJ(E,z_E)^{2}\le C_{\scFJ}\,\dD(E).
\]
\end{theorem}

\begin{proof}
\emph{Attainment.}  By Lemma~\ref{lem:Phi-cts}, \(\Phi_E\) is
continuous on \(K\), and \(K\) is compact; hence the minimum is
attained.

\emph{Upper bound at \(z_{\rm FJ}\).}  Let \(z_F,z_{\rm FJ}\) be the
Fraenkel and FJ centres of Theorems~\ref{thm:FMP}, \ref{thm:FJ}.  By
Lemma~\ref{lem:localise-centres} (after choosing \(\eta_{0}\) small),
both lie in \(K\).  By Theorem~\ref{thm:FJ},
\(\bFJ(E,z_{\rm FJ})^{2}\le C_{FJ}\,\dD(E)\).  By
Lemma~\ref{lem:change-centre} and Theorem~\ref{thm:centres-close},
\[
   \alphaFJ(E,z_{\rm FJ})
   \le \alphaFJ(E,z_F)+\frac{n\omega_{n-1}}{\omega_n}
        |z_{\rm FJ}-z_F|
   \le C_F^{1/2}\dD(E)^{1/2}
      +\frac{n\omega_{n-1}}{\omega_n}C_{cc}\dD(E)^{1/2}.
\]
Squaring,
\[
   \alphaFJ(E,z_{\rm FJ})^{2}\le C'(n)\,\dD(E).
\]
Summing,
\[
   \Phi_E(z_{\rm FJ})
   =\alphaFJ(E,z_{\rm FJ})^{2}+\bFJ(E,z_{\rm FJ})^{2}
   \le \bigl(C'(n)+C_{FJ}\bigr)\dD(E)
   =: C_{\scFJ}(n)\,\dD(E).
\]
Since \(z_{\rm FJ}\in K\),
\(m_E=\min_{K}\Phi_E\le\Phi_E(z_{\rm FJ})\le C_{\scFJ}\dD(E)\), as
claimed.
\end{proof}

\begin{remark}[On the constants]\label{rem:constants}
The constant \(C_{\scFJ}(n,K)\) does not depend on \(K\) in the proof
above; the dependence on \(K\) is only through the containment
condition \(E\subset K\) (used in
Lemma~\ref{lem:localise-centres}), which we encode by allowing
\(\eta_0\) to depend on \(K\) so that the Fraenkel and FJ centres lie
in \(K\).  Thus we may write \(C_{\scFJ}=C_{\scFJ}(n)\) once \(K\) is
fixed.
\end{remark}

\section{Borel measurable selection on torsion levels}\label{sec:selection}

We now assemble the measurable-selection input.  The setting is
\(E_{\rho}:=\{u>t(\rho)\}\subset B_R\), with \(|E_\rho|=\omega_n\rho^n\),
\(\rho\in[\rstar,1]\), and isoperimetric defect
\(\mathcal I(\rho):=D_I(t(\rho))=(P(E_\rho)-P(B_\rho))(P(E_\rho)+P(B_\rho))\).
At good (a.e.\ finite-perimeter) levels
\(P(E_\rho)-P(B_\rho)\le C(n,\rstar)\mathcal I(\rho)\).

For the joint functional we use the
\emph{scale-normalised} \(\Phi\): writing \(\hat E_\rho:=\rho^{-1}E_\rho\)
(unit volume), define
\[
   \Phi(\rho,z):=\Phi_{\hat E_\rho}(z)
   = \alphaFJ(\hat E_\rho,z)^{2}+\bFJ(\hat E_\rho,z)^{2}.
\]
Equivalently, for \(z\in K\),
\(\Phi(\rho,\rho z)=\alphaFJ(E_\rho,\rho z)^{2}+\bFJ(E_\rho,\rho z)^{2}\)
by the scaling Lemma in \texttt{WIP\_FJCenterBridge.tex}.

\begin{lemma}[Joint Borel/continuity of \(\Phi\)]\label{lem:Phi-Borel}
Assume the family \(\rho\mapsto E_\rho\) is given via the torsion
coarea: \(\rho\mapsto\mathbf 1_{E_\rho}\) is Borel from
\([\rstar,1]\) into \(L^{1}(\R^{n})\), and
\(\rho\mapsto\partial^{\ast}E_\rho\) admits a Borel
parametrisation in the sense that the pushforward measures
\((\rho\mapsto\Hn\lfloor\partial^{\ast}E_\rho)\) form a Borel field
of finite Radon measures on \(\R^{n}\) (cf.\ \cite[\S15]{Maggi2012}).
Then \((\rho,z)\mapsto\Phi(\rho,z)\) is Borel measurable on
\(\{\rho:P(E_\rho)<\infty\}\times K\), and continuous in \(z\) for
each fixed \(\rho\).
\end{lemma}

\begin{proof}
\emph{Borel in \(\rho\), continuous in \(z\).}  By
Lemma~\ref{lem:Phi-cts}, \(z\mapsto\Phi(\rho,z)\) is continuous on
\(K\) for each \(\rho\) with \(P(E_\rho)<\infty\).  For Borel
measurability in \(\rho\), note that:
\begin{enumerate}
\item \(\rho\mapsto\alphaFJ(E_\rho,\rho z)^{2}=
   \bigl(\|\mathbf 1_{E_\rho}-\mathbf 1_{B_\rho(\rho z)}\|_{L^{1}}/
      |B_\rho|\bigr)^{2}\) is Borel: \(L^{1}\)-norms of Borel
   \(L^{1}\)-valued maps are Borel, and the translation
   \(\rho\mapsto\mathbf 1_{B_\rho(\rho z)}\) is continuous.
\item \(\rho\mapsto\betaFJ(E_\rho,\rho z)^{2}=
   \tfrac12\int|\nu_{E_\rho}-(x-\rho z)/|x-\rho z||^{2}\,
   d\Hn\lfloor\partial^{\ast}E_\rho\) is Borel: it is the integral
   of a bounded continuous function on \(\R^n\setminus\{\rho z\}\)
   against a Borel field of finite Radon measures
   \cite[Prop.~15.10]{Maggi2012}.
\end{enumerate}
Sums, ratios, and squares of Borel functions are Borel, and joint
Borel measurability \((\rho,z)\mapsto\Phi(\rho,z)\) follows from
Borel-in-\(\rho\) + continuous-in-\(z\)
(see e.g.\ \cite[Lem.~III.14]{CastaingValadier1977}; alternatively
sample on a countable dense net in \(K\)).
\end{proof}

\begin{proposition}[Borel same-centre selection]\label{prop:selection}
Let \(R\ge1\), \(\rstar\in(0,1)\), and \(K:=\overline{B_{R+1}}\subset
\R^{n}\) (so \(\overline{B_{2}}\subset K\)).  Set
\[
   G_I:=\bigl\{\rho\in[\rstar,1]:
        \mathcal I(\rho)\le\eta_I\bigr\},
\]
where \(\eta_I=\eta_I(n,R,\rstar)>0\) is small enough that
Theorem~\ref{thm:main-joint} applies to \(\hat E_\rho=\rho^{-1}E_\rho\)
on \(G_I\).  Then there exist
\(C=C(n,R,\rstar)>0\) and a Borel map \(\rho\mapsto z_\rho\in K\)
defined on \([\rstar,1]\) (with the convention \(z_\rho:=0\) for
\(\rho\notin G_I\)) such that for every \(\rho\in G_I\),
\[
   \alphaFJ(\hat E_\rho,z_\rho/\rho)^{2}
      +\bFJ(\hat E_\rho,z_\rho/\rho)^{2}
      \le C_{\scFJ}\,\dD(\hat E_\rho)
   \le C\,\mathcal I(\rho),
\]
and in particular
\[
   |E_\rho\Delta B_\rho(z_\rho)|^{2}+\!\!\int_{\partial^{\ast}E_\rho}
      \!\Bigl|\nu_{E_\rho}-\tfrac{x-z_\rho}{|x-z_\rho|}\Bigr|^{2}d\Hn
   \le C\,\mathcal I(\rho).
\]
\end{proposition}

\begin{proof}
Fix \(C_0:=2C_{\scFJ}(n)\), and define the multifunction
\[
   \Gamma:G_I\rightrightarrows K,\quad
   \Gamma(\rho):=\bigl\{z\in K:
   \Phi(\rho,z/\rho)\le C_0\,\dD(\hat E_\rho)\bigr\}.
\]
By Theorem~\ref{thm:main-joint}, \(\Gamma(\rho)\ne\emptyset\) for
every \(\rho\in G_I\).  By Lemma~\ref{lem:Phi-cts} (continuity in the
centre) and the continuity of the map \(z\mapsto z/\rho\),
\(\Gamma(\rho)\) is closed in \(K\), hence compact.

\smallskip\emph{Borel measurability of the graph.}  The graph
\(\mathrm{Gr}(\Gamma)=\{(\rho,z)\in G_I\times K:\Phi(\rho,z/\rho)\le
C_0\dD(\hat E_\rho)\}\) is Borel by Lemma~\ref{lem:Phi-Borel}: the
inequality involves Borel functions of \((\rho,z)\) (and \(\rho\mapsto
\dD(\hat E_\rho)\) is Borel, being a Borel function of
\(P(E_\rho)\)).

\smallskip\emph{Kuratowski--Ryll-Nardzewski selector.}  The closed-
valued multifunction \(\Gamma\) on the Borel space \(G_I\) with
values in the Polish space \(K\) has Borel graph, hence is Borel
measurable in the sense of Kuratowski--Ryll-Nardzewski (cf.\
\cite[Theorem~III.6 and Theorem~III.30]{CastaingValadier1977}); by
the same theorem, there exists a Borel selector
\(\rho\mapsto z_\rho\in\Gamma(\rho)\).  Extending by \(z_\rho:=0\) on
\([\rstar,1]\setminus G_I\) preserves Borel measurability.

\smallskip\emph{Quantitative estimates.}  For \(\rho\in G_I\),
\(\dD(\hat E_\rho)=\dD(E_\rho)=(P(E_\rho)-P(B_\rho))/P(B_\rho)\le
C(n,\rstar)\mathcal I(\rho)\).  The joint inequality
\(\Phi_{\hat E_\rho}(z_\rho/\rho)\le C_0\,\dD(\hat E_\rho)
\le C\,\mathcal I(\rho)\) holds by definition of \(\Gamma\).
Multiplying out scales (Lemma in
\texttt{WIP\_FJCenterBridge.tex}, restated as
\(|E_\rho\Delta B_\rho(z_\rho)|=\rho^{n}|\hat E_\rho\Delta B_1(z_\rho/\rho)|\)
and similarly for the normal-oscillation integral, using
\(P(E_\rho)/P(B_\rho)\le 1+C\dD(E_\rho)\)), one obtains the two
unnormalised displayed bounds.
\end{proof}

\section{Corollary for the torsion application}\label{sec:torsion}

\begin{corollary}[Replacement for the conditional centre bridge]
\label{cor:torsion-import}
Let \(\Omega\subset B_R\) be the bounded torsion domain of the Plan~2
setting, \(E_\rho=\{u>t(\rho)\}\) its volume-radius parametrised
superlevels, \(\rstar\in(0,1)\), and \(K:=\overline{B_{R+1}}\).  For
every \(\rho\) in the good set \(G_I\) of
Proposition~\ref{prop:selection} (with \(\eta_I=\eta_I(n,R,\rstar)\)
chosen so that
Theorem~\ref{thm:main-joint} applies), the centre \(z_\rho\in K\)
selected by Proposition~\ref{prop:selection} is Borel in \(\rho\) and
satisfies the three conclusions of the
good-level centre proposition of \texttt{WIP\_FJCenterBridge.tex}
(Proposition~3.1 there, ``Good-level same-centre estimates''), with
the same constant \(C=C(n,R,\rstar)\).  In particular,
\begin{enumerate}
\item \(|z_\rho|\le R+1\),
\item \(|E_\rho\Delta B_\rho(z_\rho)|^{2}\le C\,\mathcal I(\rho)\),
\item \(\int_{\partial^{\ast}E_\rho}
   |\nu_\rho-(x-z_\rho)/|x-z_\rho||^{2}\,d\Hn\le C\,\mathcal I(\rho)\).
\end{enumerate}
\end{corollary}

\begin{proof}
The first conclusion is by construction (\(z_\rho\in K\)).  The
remaining two are the unnormalised forms of the joint estimate in
Proposition~\ref{prop:selection}, obtained by writing
\(\alphaFJ(\hat E_\rho,z_\rho/\rho)^{2}=
   (|E_\rho\Delta B_\rho(z_\rho)|/|B_\rho|)^{2}\),
multiplying by \(|B_\rho|^{2}\), and using \(\rho\ge\rstar\); and by
writing \(\bFJ(\hat E_\rho,z_\rho/\rho)^{2}\cdot P_{\hat E_\rho}^{\ast}
   =\tfrac12\rho^{-(n-1)}\int_{\partial^{\ast}E_\rho}|\nu-e|^{2}d\Hn\),
multiplying out, and using \(P(B_\rho)=n\omega_n\rho^{n-1}\),
\(P(E_\rho)\le 2P(B_\rho)\) on \(G_I\).
\end{proof}

\begin{remark}[What is unconditional, and what is cited]
\label{rem:status}
\textbf{Unconditional.}  Theorems~\ref{thm:main-joint},
\ref{thm:centres-close}, Lemma~\ref{lem:Phi-cts},
Lemma~\ref{lem:change-centre},
Lemma~\ref{lem:localise-centres}, Lemma~\ref{lem:Phi-Borel}, and
Proposition~\ref{prop:selection}.  Each is proved here from the
following \emph{cited} inputs:
\begin{itemize}
\item Theorem~\ref{thm:FMP} (sharp Fraenkel asymmetry stability):
   \cite{FigalliMaggiPratelli2010,FuscoMaggiPratelli2008}.
\item Theorem~\ref{thm:FJ} (sharp normal-oscillation stability):
   \cite{FuscoJulin2014}.
\item Fuglede's nearly-spherical lemma
   \cite{Fuglede1989} and its extension to small-deficit finite
   perimeter sets \cite[Theorem~3.4]{CicaleseLeonardi2012},
   used in Theorem~\ref{thm:centres-close} (Step~1).
\item Coarea Borel structure of reduced boundaries
   \cite[\S15]{Maggi2012}, used in Lemma~\ref{lem:Phi-Borel}.
\item Kuratowski--Ryll-Nardzewski measurable selection
   \cite[Theorem~III.6, III.30]{CastaingValadier1977}, used in
   Proposition~\ref{prop:selection}.
\end{itemize}
\textbf{Remaining external dependency.}  The convex-quadratic
inequality of Theorem~\ref{thm:centres-close}, Step~3
(Remark~\ref{rem:step3}), uses two standard ingredients of the
Fuglede/Fusco--Julin analysis (an \(L^{2}\)-orthogonality of the
sphere harmonics of degree~1 and an \(L^{\infty}\)-bound on the
nearly-spherical parametrisation); these are spelled out in
\cite[Sec.~2--3]{FuscoJulin2014} and
\cite[Sec.~3]{CicaleseLeonardi2012} but are not reproduced here.
This is the only step of the present file that is not internally
self-contained.
\end{remark}

\begin{thebibliography}{99}

\bibitem{CastaingValadier1977}
C.~Castaing, M.~Valadier,
\emph{Convex Analysis and Measurable Multifunctions},
Lecture Notes in Math.~580, Springer, 1977.

\bibitem{CicaleseLeonardi2012}
M.~Cicalese, G.P.~Leonardi,
A selection principle for the sharp quantitative isoperimetric
inequality,
Arch.\ Ration.\ Mech.\ Anal.\ \textbf{206} (2012), 617--643.

\bibitem{FigalliMaggiPratelli2010}
A.~Figalli, F.~Maggi, A.~Pratelli,
A mass transportation approach to quantitative isoperimetric
inequalities,
Invent.\ Math.\ \textbf{182} (2010), 167--211.

\bibitem{Fuglede1989}
B.~Fuglede,
Stability in the isoperimetric problem for convex or nearly spherical
domains in \(\R^{n}\),
Trans.\ Amer.\ Math.\ Soc.\ \textbf{314} (1989), 619--638.

\bibitem{FuscoJulin2014}
N.~Fusco, V.~Julin,
A strong form of the quantitative isoperimetric inequality,
Calc.\ Var.\ Partial Differential Equations \textbf{50} (2014),
925--937.

\bibitem{FuscoMaggiPratelli2008}
N.~Fusco, F.~Maggi, A.~Pratelli,
The sharp quantitative isoperimetric inequality,
Ann.\ of Math.\ \textbf{168} (2008), 941--980.

\bibitem{Maggi2012}
F.~Maggi,
\emph{Sets of Finite Perimeter and Geometric Variational Problems},
Cambridge Stud.\ Adv.\ Math.~135, Cambridge Univ.~Press, 2012.

\end{thebibliography}

\end{document}
